% ch-a4-15.tex —— 第 4 篇（Emerton–Gee–Hellmann）续片，底本 PDF 348–359 = 印刷页 335–346

\ysrc{335}
……（如期望 7.6.12 所预言）的方式如下：$\overline{\rho}$ 的 Serre 权 $\sigma$ 是已知的，它们恰是 $M^\vee_\infty[\mathfrak m_\infty]$ 的 $K$-基座（$K$-socle）的不可约成分，且已知它们以重数一出现。这些权如何分布到 $M^\vee_\infty[\mathfrak m_\infty]$ 的各个直和项之中，也是已知的。（这些结果见 [GLS14, GK14, EGS15, Bre14]。）于是，当 $\sigma$ 遍历 $\overline{\rho}$ 的 Serre 权时，重数空间 $\Hom_G(c\text{-}\Ind_{KZ}^G\sigma,M^\vee_\infty[\mathfrak m_\infty])$ 是一维的，且已知任何非零态射经过 $M^\vee_\infty[\mathfrak m_\infty]$ 的哪个直和项分解。

另一方面，我们可以把这个 Hom 空间写成
\[
\begin{aligned}
\Hom_G\bigl(c\text{-}\Ind_{KZ}^G\sigma,M^\vee_\infty[\mathfrak m_\infty]\bigr)
&=\Hom_G\Bigl(c\text{-}\Ind_{KZ}^G\sigma,\mathfrak B(\Delta_x)\Bigr)\\
&=\Hom_G\Bigl(c\text{-}\Ind_{KZ}^G\sigma,\mathfrak B\bigl(\oplus_n\delta_x(n)\bigr)\Bigr)\\
&=\Hom_G\Bigl(c\text{-}\Ind_{KZ}^G\sigma,\oplus_n\mathfrak B\bigl(\delta_x(n)\bigr)\Bigr)\\
&=\oplus_n\Hom_G\Bigl(c\text{-}\Ind_{KZ}^G\sigma,\mathfrak B\bigl(\delta_x(n)\bigr)\Bigr),
\end{aligned}
\]
其中求和遍历 $n\in\Z$ 或 $n\in\Z/2\Z$，视 $\overline{\rho}$ 是特征标的直和还是不可约而定。现在，$\mathfrak A$ 与 $\mathfrak B$ 之间的伴随性表明
\[
\begin{aligned}
\Hom_G\Bigl(c\text{-}\Ind_{KZ}^G\sigma,\mathfrak B\bigl(\delta_x(n)\bigr)\Bigr)
&=\Hom_G\Bigl(\mathfrak A\bigl(c\text{-}\Ind_{KZ}^G\sigma\bigr),\delta_x(n)\Bigr)\\
&=\Hom_G\Bigl(\OO_{\XX(\sigma)}(-1,-1),\delta_x(n)\Bigr)\\
&=i^*\bigl(\OO_{\XX(\sigma)}(-1,-1)\bigr)(n)^{G_x},
\end{aligned}
\]
其中如前所述，$i:[\Spec k/G_x]\into\XX$ 是 $x$ 处的剩余带叠（residual gerbe）到 $\XX$ 的浸入。由于 $\OO_{\XX(\sigma)}(-1,-1)$ 是可逆层，我们看到 $i^*\bigl(\OO_{\XX(\sigma)}(-1,-1)\bigr)$ 是 $G_x$ 的一个特征标，因而形如 $k(m)$，其中 $m$ 是某个整数（我们看到它由 $\sigma$ 与 $x$ 决定）。于是 $i^*\bigl(\OO_{\XX(\sigma)}(-1,-1)\bigr)(n)^{G_x}=k(m+n)^{G_x}$ 恰对一个 $n$ 值非零（即 $n=-m$），而这就确定了 $M^\vee_\infty[\mathfrak m_\infty]$ 的相应成分的扭；也就是说，该成分以 $m$ 作扭而出现。对 $x$（也就是对 $\overline{\rho}$）与 $\sigma$ 的各种可能选择计算 $m$ 的值，便得到例 (4) 与 (5) 中出现的那些扭。

\end{rem}

\begin{rem}{7.8.5}\label{rem:a4-7-8-5}
正如讨论 §7.7.59 时已经回顾过的，[Sch15b] 证明了 $\mathrm{GL}_2(\Q_{p^2})$ 的超奇异表示从不是有限表现的。把这一事实与前述（猜想性的）计算相比较，我们看到不应指望伴随函子 $\mathfrak B$ 保持紧对象。（我们已经看到，它应当把某些摩天楼层——这些层是凝聚的，因而作为 $\XX$ 上 Ind-凝聚层构成的稳定 $\infty$-范畴的对象是紧的——送到 $\mathrm{GL}_2(\Q_{p^2})$ 的非有限表现表示，而后者作为 $\Ind\,D^b_{\mathrm{f.p.}}(\mathrm{sm}\,G)$ 的对象不是紧的。）

要看出其中缘由，注意伴随函子 $\mathfrak B$ 应当可以用 $\Ind\,\Coh(\XX)$ 的一个半正交分解来计算，该分解正是 §7.5.11 所猜想的那一类。确切地说，猜想中的函子 $\mathfrak A$ 应当诱导 $\Ind\,D^b_{\mathrm{f.p.}}(\mathrm{sm}\,G)$ 与 $\Ind\,\Coh(\XX)$ 的某个半正交分解的「第一块」之间的等价，而该分解的其余各块应当由「几何 Eisenstein 级数」型的构造得到，这些构造涉及从 $\XX$ 中可约轨迹出发的推前（以及来自 $\mathcal{BC}(\OO(1))$ 与 $\mathcal{BC}(\OO(2))$ 的更为神秘的贡献）。复合函子 $\mathfrak A\circ\mathfrak B$ 于是可以描述为「向远离」这些其余各块的方向作投影，以落入 $\mathfrak A$ 的本质像（从而 $\mathfrak B$\ysrc{336} 本身可以描述为这个投影，再对本质像中所得的对象施加 $\mathfrak A^{-1}$）。
\end{rem}

由于典型不可约分支 $\XX(\sigma)$ 的可约轨迹到 $\XX(\sigma)$ 中的包含是开浸入而非闭浸入（因而非固有），半正交分解中这些「其余各块」里的对象通常不凝聚。于是，对凝聚层求 $\mathfrak A\circ\mathfrak B$ 的值，通常涉及取从这些非凝聚层到给定凝聚层的态射的锥，所以我们一般不应指望 $\mathfrak A\circ\mathfrak B$ 在凝聚层上的取值自身还是凝聚的。相应地，我们一般也不应指望 $\mathfrak B$ 在凝聚层上的取值是 $\Ind\,D_{\mathrm{f.p.}}(\mathrm{sm}\,G)$ 的紧对象。

\begin{rem}{7.8.6}\label{rem:a4-7-8-6}
自 Breuil 与 Paškūnas 的结果 [BP12] 以来，已经清楚：$\mathrm{GL}_2(\Q_{p^2})$ 的 $p$-进局部 Langlands 对应的表述必定与 $\mathrm{GL}_2(\Q_p)$ 的（经典）对应大不相同，因为 $\mathrm{GL}_2(\Q_{p^2})$ 的表示似乎远比 $\Gal_{\Q_{p^2}}$ 的二维表示多。特别地，正如讨论 §7.7.59 时已经讨论过的，$\mathrm{GL}_2(\Q_{p^2})$ 有无限族的超奇异不可约容许 mod $p$ 表示，它们似乎都对应于 $\Gal_{\Q_{p^2}}$ 的同一个不可约表示 $\overline{\rho}$。

我们猜想中的函子 $\mathfrak A$ 为这一现象提供了一种解释：伽罗瓦表示的表面匮乏，由 $\XX$ 上凝聚层的丰富所补偿；特别地，这些超奇异表示都对应于 $\XX$ 上的 Ind-凝聚层，而这些层由其（拟凝聚但非凝聚的）$H^{-1}$ 层区分开来。

不过，还有一种更早的途径来调和 [BP12] 的结果与「$\mathrm{GL}_2(\Q_{p^2})$ 的 $p$-进局部 Langlands 对应应当存在」这一预期：人们可以希望，在 [BP12] 所构造的无限族超奇异表示之中，存在唯一确定的典范元素 $\pi_{\overline{\rho}}$。而且，这个表示 $\pi_{\overline{\rho}}$ 应当由局部--整体相容性刻画：对 $\overline{\rho}$ 的任何整体化，它都应当出现在志村曲线的上同调中。在这个意义上，$\pi_{\overline{\rho}}$ 的存在相当于证明：志村曲线上同调中出现的表示只依赖于 $\overline{\rho}$——并且这个表示确实是 [BP12] 所构造类型的那种超奇异不可约表示。许多作者研究过这一问题，支持这一看法的证据相当多；我们再次请读者参看 [BHH+21] 的引言，那里对这些研究作了简明总结。

我们坚信典范表示 $\pi_{\overline{\rho}}$ 存在。事实上，正是这一视角促成了 [CEG+16] 中经由 Taylor--Wiles 拼接来构造 $p$-进局部 Langlands 对应的方案，而该方案又启发了本讲义中的诸猜想（见注 6.1.31）。我们将在下文 §9.3 中更完整地讨论局部--整体相容性。本节的这些计算也与该预期相容；确实，由上文例 (4) 与 (5) 我们看到，若 $x$ 是对应于适当一般（generic）的不可约 $\overline{\rho}$ 的闭点，则应当有 $\pi_{\overline{\rho}}=\mathfrak B(\delta_x)$。

此外，在我们的猜想框架内，我们断言：可以通过下述性质在所有 $\mathrm{GL}_2(\Q_{p^2})$ 的不可约表示 $\pi$ 中刻画 $\pi_{\overline{\rho}}$：它是使 $H^0(\mathfrak A(\pi))$ 的支撑包含 $x$ 的唯一不可约 $\pi$。事实上，设 $\pi$ 不可约，且 $H^0(\mathfrak A(\pi))$ 的支撑

\ysrc{337}
包含 $x$。由于 $\pi$ 与 $\pi_{\overline{\rho}}=\mathfrak B(\delta_x)$ 都不可约，只需证明
\[
\Hom_{\mathrm{GL}_2(\Q_{p^2})}\bigl(\pi,\mathfrak B(\delta_x)\bigr)\neq0.
\]
由伴随性，这等价于证明
\[
\Hom_{\mathrm{Pro}\text{-}\OO_\XX}\bigl(\mathfrak A(\pi),\delta_x\bigr)\neq0.
\]
为看出这一点成立，我们考虑典范态射
\[
\mathfrak A(\pi)\to\tau^{\ge0}\mathfrak A(\pi),
\]
并注意 $\tau^{\ge0}\mathfrak A(\pi)$ 凝聚（由注 6.1.29），故由「$H^0(\mathfrak A(\pi))$ 的支撑包含 $x$」这一假设，它到 $\delta_x$ 有态射（这里我们用到了 $\mathfrak A(\pi)$ 落在 $\Ind\,\Coh(\XX)$ 的偶部之中，见 §7.5.11）。

注意，由于超奇异表示 $\pi_{\overline{\rho}}$ 可从一个表示 $c\text{-}\Ind_{KZ}^G\sigma$（其中 $\sigma$ 是某个 Serre 权）得到满射，凝聚层 $\tau^{\ge0}\mathfrak A(\pi)$ 可从一个线丛得到满射。特别地，在情形 (5) 中我们看到 $H^0(\mathfrak A(\pi))$ 应当等于 $\delta_x$。
\end{rem}

\begin{rem}{7.8.7}\label{rem:a4-7-8-7}
在上面的情形 (4) 与 (5) 中，我们没有 $\mathfrak A(\pi)$ 的显式描述，只有情形 (5) 中它的 $H^0$ 等于 $\delta_x$（如注 7.8.6 所讨论）。（类似地，我们预期在情形 (4) 中 $H^0$ 是所有半单化为 $x$ 的点的闭轨迹的结构层。）不过，$\mathfrak A(\pi)$ 原则上可以用 §A.8.12 的机器来计算，因为它等于 $(\mathfrak A\circ\mathfrak B)(\delta_x)$；并且我们预期能够写出猜想函子 $\mathfrak A$ 的像的一组生成元（即对应于各个 $\mathfrak A(c\text{-}\Ind_{KZ}^G\sigma)$ 的凝聚层，在一般情形下我们已在 §7.7 中显式写出）。
\end{rem}

\begin{rem}{7.8.8}\label{rem:a4-7-8-8}
对 $\mathrm{GL}_2(F)$ 的传统 mod $p$ 局部 Langlands 对应，一个预期中的特征是：对可约的 $\overline{\rho}$，若（沿用上面引入的记号）记 $\pi(\overline{\rho}):=\mathfrak B(\Delta_x)$，则应当有
\begin{equation}\tag{7.8.9}\label{eq:a4-7-8-9}
\pi(\overline{\rho}^{\mathrm{ss}})\cong\pi(\overline{\rho})^{\mathrm{ss}}.
\end{equation}
自然要问：从我们当前的视角能否看出这一点。

为此，我们注意：若 $\overline{\rho}^{\mathrm{ss}}$ 是 $G_{\Q_{p^f}}$ 的两个适当一般的特征标 $\chi_1$ 与 $\chi_2$ 的直和\tn{底本此处作 $G_{\Q_{p^f}}$，与他处 $\Gal_{\Q_{p^f}}$ 的记号不一致，据上下文当指该绝对伽罗瓦群；译文照录原文（上下文推断，无外部依据）。}，则作为 $\chi_2$ 被 $\chi_1$ 扩张的 $\overline{\rho}$ 的轨迹等于 $[\mathbf A^f/\mathbf G_m]$ 的一个拷贝。若以 $\mathcal F$ 记 $[\mathbf A^f/\mathbf G_m]$ 的结构层推到 $\XX$ 上所得的凝聚层，则 $\mathcal Hom(L_\infty,\mathcal F)$ 可以解释为 $\mathbf A^f$ 上光滑 $\mathrm{GL}_2(F)$-表示的一个 $\mathbf G_m$-等变平坦族。若忘记 $\mathbf G_m$-等变结构，这个族就在由 $\mathbf A^f$ 参数化的那些伽罗瓦表示（即特征标的扩张）$\overline{\rho}$ 上把这些 $\pi(\overline{\rho})$ 插值起来。看来可以相信，这个平坦族（连同其 $\mathbf G_m$-等变结构！）的存在能够蕴含所预期的同构 (7.8.9)。
\end{rem}

\subsection{全局函数与 Hecke 代数}
猜想函子 $\mathfrak A$ 的全忠实性对叠 $\XX_d$ 的几何有一些具体的推论。例如，我们在上文看到，若 $G=\mathrm{GL}_2(F)$ 且 $F/\Q_p$ 非分歧，则对一般的 Serre 权 $\sigma$，我们预期 $\mathfrak A(c\text{-}\Ind_K^G\sigma)$ 是分支 $\XX(\sigma)$ 上的线丛；特别地，（非导出的）Hecke 代数 $\End_G(c\text{-}\Ind_K^G\sigma)$ 应当等同于 $\XX(\sigma)$ 上的全局函数。

\ysrc{338}
这一预期直接由 $\XX(\sigma)$ 借助 Fontaine--Laffaille 理论的显式描述得到；对驯顺型，用显式的整 $p$-进 Hodge 理论也可以证明类似结果（同样只对 $F/\Q_p$ 非分歧的 $\mathrm{GL}_2(F)$ 的一般情形）。

更一般地，下述结果由 Eivind Otto Hjelle、Louis Jaburi、Rachel Knak、Hao Lee 与 Shenrong Wang 的论文 [HJK+24] 证明：若 $\sigma$ 是 $\mathrm{GL}_d(F)$ 的一个适当一般的 Serre 权（$F/\Q_p$ 非分歧），且 $\XX(\sigma)$ 是叠 $\XX_{F,d}$ 中由 $\sigma$ 标号的不可约分支，则 $\XX(\sigma)$ 上的全局函数由形如 $k[T_1,\dots,T_{d-1},T_d,T_d^{-1}]$ 的环给出，而该环同构于 Hecke 代数 $\End_G(c\text{-}\Ind_K^G\sigma)$。尽管 [HJK+24] 并未在这个 Hecke 代数与该全局函数环之间定义具体的映射，这样的映射随后已由 Heejong Lee 构造出来 [Lee24]，他在其中将它解释为 [Her11] 的 mod $p$ Satake 同构的几何类似物。

在非一般的情形，Anthony Guzman、Kalyani Kansal、Jason Kountouridis、Ben Savoie 与 Xiyuan Wang 的论文 [GKK+25] 在关于 $\sigma$ 的温和假设下证明：对 $F/\Q_p$ 非分歧、$\sigma$ 是 $\mathrm{GL}_2(F)$ 的一个未必一般的 Serre 权，$\XX(\sigma)$ 的正规化上的全局函数可以等同于相应的 Hecke 代数。这些结果又由 Kansal 与 Savoie 在其论文 [KS25] 中进一步精化。

\section{$\ell\neq p$ 情形的范畴局部 Langlands}

尽管我们迄今的焦点一直放在 $p$-进局部域上群 $\mathrm{GL}_d$ 的 $p$-进局部 Langlands 猜想的范畴表述，我们也希望讨论整体情形；确实可以这么说：发展数论的局部方面，其主要目的正是为研究整体问题提供工具。我们将在 §9 中解释整体理论；这自然涉及局部的 $\ell\neq p$ 情形，也就是本节所研究的情形。

我们将聚焦的主要整体问题，是用范畴局部猜想所产生的层来猜想性地描述志村簇（更带推测性地说，一般同余商）的上同调。一方面，这与 Fontaine--Mazur 猜想的完备上同调途径有关（如 [Eme11a, Pan22a]）。另一方面，它与 Taylor--Wiles 拼接有关（特别地，它为拼接模与范畴猜想中凝聚层之间预期中的关系提供了一种猜想性的解释）。

为研究整体问题，特别是涉及志村簇的问题，我们必须能够考虑未必分裂的约化群。尽管目前我们对非 $\mathrm{GL}_d$ 情形的 $p$-进局部 Langlands 的理解比对 $\mathrm{GL}_d$ 的理解还要零碎得多，但在远离 $p$ 的素位处，局部情形被理解得更好；确实，此时有精确的猜想，并且已经取得了可观的进展，特别是在论文 [FS24, Zhu25a, Hel23, BZCHN24, Zhu25b] 中。本节我们简要回顾其中一些局部猜想与结果，细节请参看这些论文以及笔记 [FS25]。

我们首先简要回顾 C-群的形式体系；当涉及非分裂群时，C-群将是局部语境中出现的伽罗瓦表示的目标。然后我们描述

\ysrc{339}
$\ell\neq p$ 情形下局部 Langlands 参数的模叠，并非常简略地回顾为整体讨论所需的范畴局部 Langlands 猜想。

\subsection{C-群}
一如既往，我们固定了 $\Q_p$ 的一个有限扩张 $L$，其整数环 $\OO$ 充当我们的系数环。

若 $F$ 是 $\Q_\ell$（$\ell$ 为某个素数）或 $\Q$ 的有限扩张，$G$ 是 $F$ 上的连通约化线性代数群，则我们可以按 [BG14] 构造 $G$ 的 C-群 ${}^cG$；另见 [Zhu25a, §3.0]。这是 L-群 ${}^LG$ 的一个变体，更适合处理与定义域相关的问题。按构造，它是半直积
\[
{}^cG:=\hat G\rtimes\bigl(\GG_m\times\Gal(\tilde F/F)\bigr),
\]
其中 $\tilde F$ 是 $F$ 的某个有限伽罗瓦扩张，按当前需要选得「足够大」，$\hat G$ 是通常的 Langlands 对偶群。（特别地，$\tilde F$ 应包含 $G$ 的拟分裂内型的分裂域；下文出现的涉及 ${}^cG$ 的定义容易看出与 $\tilde F$ 的具体选择无关。）这里 $\GG_m$ 经由余特征标 $\rho_{\mathrm{ad}}:\GG_m\to\hat G_{\mathrm{ad}}$ 作用，该余特征标由 $\hat G$ 的正余根之和的一半定义。我们把 ${}^cG$ 看作整数环 $\OO$ 上的群概形。

\begin{rem}{8.1.1}\label{rem:a4-8-1-1}
C-群也可以定义为 $G$ 的某个 $\GG_m$-扩张 $\tilde G$ 的 L-群，见 [BG14, Defn. 5.3.2]。
\end{rem}

\subsection{局部情形下的 Langlands 参数：$\ell\neq p$}
若 $F$ 是 $\ell$-进域且 $\ell\neq p$，则可以定义一个叠 $\XX_{G,F}$，它参数化 $F$ 的 Weil--Deligne 型群（作适当解释）到 ${}^cG$ 的表示。事实上，定义 $\XX_{G,F}$ 有各种方式，取决于我们希望它生活在什么基之上。我们把 $\XX_{G,F}$ 看作 $\Spf\OO$ 上的叠，于是可以作如下定义：先选取元素 $\sigma,\tau\in\Gal_F:=\Gal(\overline{F}/F)$，使它们分别提升 Frobenius 与驯顺惯性的一个拓扑生成元，并以 $\Gamma_F$ 记 $\Gal_F$ 中由 $P_F$（$\Gal_F$ 的野惯性子群）、$\sigma$ 与 $\tau$ 生成的子群。我们有群的短正合列
\[
1\to P_F\to\Gamma_F\to\Gamma^{\mathrm{tame}}_F:=\langle\sigma,\tau\ \mid\ \sigma\tau\sigma^{-1}=\tau^q\rangle\to1
\]
（其中 $q$ 记 $F$ 的剩余域的大小）。注意 $P_F$ 有一组由 $\Gamma_F$ 的正规开子群 $Q$ 构成的单位元邻域基，于是我们可以把 $\Gamma_F=\varprojlim_Q\Gamma_F/Q$ 写成离散商的投射极限，其中每个 $Q$ 都是 $P_F$ 的子群（从而是 pro-$\ell$ 的）。

对任何使 $p$ 幂零的 $\OO$-代数 $A$，以及前面提到的任何子群 $Q$，我们定义 $\XX^Q_{G,F}(A)$ 为 C-同态 $\Gamma_F/Q\to{}^cG(A)$ 构成的群胚，即满足通常「L-同态」条件之 C-群版本的同态：复合
\begin{equation}\tag{8.2.1}\label{eq:a4-8-2-1}
\Gamma_F\to\Gamma_F/Q\to{}^cG(A)\to A^\times\times\Gal(\tilde F/F)
\end{equation}
（第二个箭头为自然的投影）应当等于逆分圆特征 $\Gamma_F\to\Gal_F\to\Z_p^\times$ 与典范满射 $\Gamma_F\to\Gal_F\to\Gal(\tilde F/F)$ 之积。然后置 $\XX_{G,F}:=\varinjlim_Q\XX^Q_{G,F}$；每个转移映射都是开闭浸入（因为每个 $Q$ 都是 pro-$\ell$ 的且 $\ell\neq p$），于是 $\XX_{G,F}$ 被赋予 Ind-代数叠的结构。

\ysrc{340}
其实，上面的讨论把 $\XX_{G,F}$ 定义为一个经典叠，但我们真正想把它定义为导出叠，这可以按附录 F 所解释的想法来做。关键的一点是：即使把 $\XX_{G,F}$ 定义为导出叠，结果它仍是经典的，甚至是局部完全交（[Zhu25a]，命题 3.1.6）。由于局部伽罗瓦上同调支撑在上同调次数 $[0,2]$ 中，这相当于说 $\XX_{G,F}$ 底层的经典叠具有预期的维数；这一点可以用显式计算证明，下面我们简要回顾。

第一点是：对 $\XX_{G,F}$ 的一个 $A$-值点 $\rho:\Gamma_F\to{}^cG(A)$（这里 $A$ 是使 $p$ 幂零的（经典！）$\OO$-代数），$\XX_{G,F}$ 在 $\rho$ 处的切复形由平移后的连续上链复形 $C^\bullet(\Gamma_F,\mathfrak g\otimes_\OO A)[1]$ 计算，其中 $\Gamma_F$ 经由 $\rho$ 与 ${}^cG$ 在其李代数上的伴随作用的复合作用在 $\mathfrak g\otimes_\OO A$ 上。该复形的上同调支撑在次数 $[-1,1]$ 中，于是我们发现 $\XX_{G,F}$ 在任意点处的余切复形的上同调在次数 $\le-2$ 消没。这意味着 $\XX_{G,F}$ 是拟光滑的（quasi-smooth），即一个（可能导出的）局部完全交。此外，$\XX_{G,F}$ 具有其预期维数（$C^\bullet(\Gamma_F,\mathfrak g\otimes_\OO A)[1]$ 的欧拉示性数等于零，本质上由局部欧拉示性数公式得到，而 $\XX_{G,F}$ 在 $\Spf\OO$ 上平坦、相对维数为零——见 [Zhu25a]，命题 3.2.7）。这意味着 $\XX_{G,F}$ 实际上是经典局部完全交，如所断言（见 [Zhu25a]，命题 2.2.13 的证明；[Zhu25a]，命题 3.1.6 的证明最终依赖于它）。

\subsection{范畴局部 Langlands}
设 $F$ 是 $\ell$-进局部域，$\ell\neq p$ 为某个素数，$G$ 是 $F$ 上的连通约化群。以 $D(G(F))$ 记 $G(F)$ 在挠 $\OO$-模上的光滑表示构成的复形的稳定 $\infty$-范畴。以 $D(G(F))'$ 记 $D(G(F))$ 的一个「重正规化」，类似于 [Zhu25a, §4.1] 中所描述的那一个（其中 $\Rep(G(F),\OO)$ 被换成 $\Rep^{\mathrm{ren}}(G(F),\OO)$，该语境与我们这里的关键差别在于那里不加挠假设）；在 $D(G(F))'$ 中，一切紧诱导表示都是紧对象（按定义）。

于是我们在 $\ell\neq p$ 语境下有如下范畴形式的局部 Langlands 猜想。

\begin{conj}{8.3.1}\label{conj:a4-8-3-1}
存在稳定 $\infty$-范畴之间的全忠实连续函子
\[
\mathfrak A_G:D(G(F))'\to\Ind\,\Coh(\XX_{G,F}).
\]
\end{conj}

\begin{rem}{8.3.2}\label{rem:a4-8-3-2}
如同注 6.1.24，我们预期猜想 8.3.1 中的函子 $\mathfrak A_G$ 应当由与某个核作（导出）张量积给出。在 $G=\mathrm{GL}_n$ 的情形，这一预期首先由 E.H. 在 [Hel23] 中明确写出；此处对 $\mathrm{GL}_n(F)$ 的 $L$-表示的 Iwahori 块，预期该函子由在 Iwahori Hecke 代数上与 Emerton--Helm--Moss 族式局部 Langlands 对应中的 Iwahori 不变向量作导出张量积给出。

猜想 8.3.1 如上所述相当不精确。Zhu 在 [Zhu25a]，猜想 4.5.1 中给出了一个精确得多的版本。（在该陈述中也没有对挠表示的限定，而目标范畴是某个代数 $\OO$-叠上的 Ind-凝聚层，我们的叠 $\XX_{G,F}$ 是它的 $p$-进完备化。）在 $L$-系数的情形，类似的猜想已\ysrc{341}由 E.H. [Hel23] 独立提出。

当 $G$ 分裂时，类似的猜想（同样在 $L$-系数情形）又由 Ben-Zvi--Chen--Helm--Nadler [BZCHN24] 独立提出，他们还构造了猜想函子在 Iwahori 块上的限制。这一限制（以及它到幂幺参数一般情形的推广）也由 Hemo--Zhu [Zhu25b] 独立构造。我们指出，当 $G$ 不形如 $\mathrm{GL}_d$（对某个 $d$）时，函子 $\mathfrak A_G$ 依赖于辅助数据的选择（例如当 $G$ 拟分裂时，一个 pinning 以及相伴的 Whittaker 数据）。叠 $\XX_{G,F}$ 上应当有线丛，使人可以在由不同辅助数据选择所产生的函子之间作扭（见 [Zhu25a]，注 4.5.7）。

在 $G=\mathrm{GL}_d$ 的情形，[BZCHN24] 的构造可推广到光滑 $\mathrm{GL}_d(F)$-表示的整个范畴（本质上是因为 $\mathrm{GL}_d$ 的 Bernstein 中心理论中出现的 Hecke 代数都是 Iwahori Hecke 代数）。正如注 8.3.2 所解释的，E.H. 以「族式局部 Langlands」给出了这个函子的另一种猜想性描述（仍用 $L$-系数），其存在性曾在 [EH14] 中猜想，并由 Helm 与 Moss [HM18] 证明。看来 E.H. 的猜想性描述应当也适用于 $\OO$-系数的情形；例如见 [Zhu25a]，注 4.4.4，那里 Zhu 在 $G=\mathrm{GL}_d$ 的情形还给出了该函子的又一种猜想性描述，其灵感来自 Braverman--Finkelberg [BF22] 在几何 Langlands 语境中所作的猜想。（据我们所知，即便在 $L$-系数的情形，也尚未证明 [BZCHN24] 的构造与 E.H. 经由族式局部 Langlands 的猜想性描述、或 Zhu 的 Braverman--Finkelberg 式猜想性描述中的任何一个相符。）

事实上，猜想 8.3.1 的全忠实函子应当实际来自几何 Langlands 对应在其范畴形式的算术类似物所给出的范畴等价的限制。Fargues 与 Scholze [FS24] 提出过这样的猜想。Zhu 也提出过另一个这样的猜想 [Zhu25a, §4.6，特别是猜想 4.6.4 与注 4.6.7]。
\end{rem}

\begin{rem}{8.3.3}\label{rem:a4-8-3-3}
设 $G=\mathrm{GL}_d(F)$。在特征零情形（在固定的 $p$-进基域 $L$ 上），猜想 8.3.1 应当有一个版本，涉及 $L$-向量空间上的光滑 $G$-表示，以及 §5.2 的叠
\[
\mathcal Y_{G,F}=\bigcup_{F'}\Mod_{d,\varphi,N,F'/F}
\]
上的（ind-）凝聚层，而非 $L$-参数的叠 $\XX_{G,F}$。注意 $\mathrm{GL}_d$ 的 $L$-参数叠与 $d$ 维 Weil--Deligne 表示的模叠 $\mathrm{WD}_{d,F}$ 一致。要点在于，我们预期「族式局部 Langlands」在 $\mathcal Y_{G,F}$ 上说得通。事实上，沿用 [Hel16]，借助 Whittaker 数据 $(N,\psi)$ 的选择以及从 $\XX_{G,F}$ 上不变函数环到 $L$-向量空间上光滑 $G$-表示范畴的 Bernstein 中心 $Z$ 的同构，可以构造 $\XX_{G,F}$ 上的一族表示 $\mathcal V$。于是 $\mathcal V$ 是
\[
(c\text{-}\Ind_N^G\psi)\otimes_{\ZZ}\OO_{\XX_{G,F}}
\]
的唯一商，使它在 $\XX_{G,F}$ 的一般点处具有某个指定的商。利用命题 5.2.10 与命题 5.2.12，同样的构造应当可以搬到 $\mathcal Y_{G,F}$ 上。这一策略有两处困难：我们需要固定一个 Whittaker 数据 $(N,\psi)$，而这只有在 $L$ 包含全部 $p$ 幂单位根时才可能（从而

\ysrc{342}
特别地 $L$ 不能是 $\Q_p$ 的有限扩张）。此外，所涉及的叠不是几何不可约的（其连通分支也不是几何连通的）。因此需要证明某些有理性结果，才能在给定的基域 $L$ 上定义族 $\mathcal V$。
\end{rem}

\section{整体 Langlands 参数的模叠与志村簇的上同调}

本节我们简要描述该理论的整体方面。我们勾勒整体伽罗瓦表示的模叠以及相伴的整体到局部限制映射的构造，然后解释如何用这些映射连同猜想性的范畴局部 Langlands 对应，给出志村簇上同调的一个猜想性公式。接着我们解释 [EGZ] 在 $\mathrm{GL}_2/\Q$ 情形得到的一些结果，这些结果把这些猜想与 Taylor--Wiles 方法以及 M.E. 关于 $p$-进 Langlands 纲领中局部--整体相容性的结果 [Eme11a] 联系起来。

\subsection{整体模叠}
我们从关于整体 Langlands 参数模叠的一些基本事实开始。虽然通常可以（通过应用限制数量）化归到考虑定义在 $\Q$ 上的约化群，但考虑一般数域的情形并不更难，我们就这么做。于是固定一个数域 $F$ 与定义在 $F$ 上的连通约化群 $G$。还固定 $F$ 的一个有限素位集 $S$，它包含 $p$ 之上与 $\infty$ 之上的全部素位。

我们可以定义一个位于 $\Spf\OO$ 之上的导出叠 $\XX_{G,F,S}$，其 $A$-值点的群胚 $\XX_{G,F,S}(A)$（对使 $p$ 幂零的经典 $\OO$-代数 $A$）与连续表示 $\rho:\Gal_{F,S}\to{}^cG(A)$ 的群胚一致（目标取离散拓扑），这些表示是 C-同态（即满足 (8.2.1)）。（这定义了 $\XX_{G,F,S}$ 在经典环上的点，也就是说它刻画了底层的经典叠 $\XX^{\mathrm{cl}}_{G,F,S}$；$\XX_{G,F,S}$ 在导出环上的点的定义更为复杂——一般形式体系见 [Zhu25a, §2.4]，本语境中具体情形的总结见 [Zhu25a, §3.4] 末尾的讨论，精确细节见即将发表的文章 [EZ]。）

当 $G=\mathrm{GL}_d$（$d\ge1$）时，底层经典叠 $\XX^{\mathrm{cl}}_{\mathrm{GL}_d,F,S}$ 已由 Wang-Erickson [WE18] 引入。他证明它是形式代数叠的不交并，其中每个都在某个完全 Noether 局部环（实际上是一个适当的伪形变环）的形式 $\Spec$ 上有限型。一般而言 $\XX_{G,F,S}$ 具有类似的结构，特别是它是形式导出代数叠。更精确地说，若 $\theta$ 是 $\Gal_{F,S}$ 的剩余伪表示，则以 $X_\theta$ 记 $\theta$ 的相伴形式形变空间。从伽罗瓦表示过渡到其相伴伪表示给出态射 $\XX_{G,F,S}\to\coprod_\theta X_\theta$；我们以 $\XX_\theta$ 记 $X_\theta$ 的原像。于是得到分解
\begin{equation}\tag{9.1.1}\label{eq:a4-9-1-1}
\XX_{G,F,S}:=\coprod_\theta\XX_\theta,
\end{equation}
它恰是 $\XX_{G,F,S}$ 按其连通分支的分解。

正如局部情形，若 $\rho:\Gal_{F,S}\to{}^cG(A)$ 是 $\XX_{G,F,S}$ 的一个 $A$-值点（其中 $A$ 是使 $p$ 幂零的经典 $\OO$-代数），则 $\XX_{G,F,S}$ 在 $\rho$ 处的切复形经由群上同调计算，此处是连续上链复形

\ysrc{343}
$C^\bullet(\Gal_{F,S},\mathfrak g)[1]$。从现在起我们假定 $p>2$。于是，正如局部情形，$\Gal_{F,S}$ 的上同调在次数 $>2$ 消没，因此我们发现 $\XX_{G,F,S}$ 在 $\Spf\OO$ 上形式拟光滑。然而与局部情形不同，$\XX_{G,F,S}$ 未必具有其预期维数（某点处的预期维数可以用 $C^\bullet(\Gal_{F,S},\mathfrak g)[1]$ 的欧拉示性数计算，后者由整体欧拉示性数公式给出，因而依赖于所考虑的点对应的伽罗瓦表示在 $F$ 的实素位处的局部行为）；更精确地说，它可以有维数大于预期维数的分支。不过，由于它形式拟光滑，它是导出局部完全交，即可以写成 ${}^c\widehat G$（${}^cG$ 的 $p$-进完备化）按适当定义的共轭作用在形如
\[
\Spf\OO[[X_1,\dots,X_d]]/(f_1,\dots,f_r)
\]
的导出形式仿射概形上的商，其中 $1+d-r-\dim{}^cG$ 等于 $\XX_{G,F,S}$ 的预期维数。（这里第一个加项「$1$」对应于 $\OO$ 的维数。）不过，这些关系 $f_1,\dots,f_r$ 未必构成正则序列，因此商环必须在导出的意义下理解。于是 $\XX_{G,F,S}$ 的导出结构支撑在 $\XX^{\mathrm{cl}}_{G,F,S}$ 那些维数大于预期维数的分支上。

\subsection{整体到局部的映射}
伽罗瓦表示的限制应当给出态射
\begin{equation}\tag{9.2.1}\label{eq:a4-9-2-1}
f:\XX_{G,F,S}\to\prod_{v\in S}\XX_{G,F_v}.
\end{equation}
我们只在 $G=\mathrm{GL}_d$ 于 $v\mid p$ 的素位处局部成立时才定义这个态射。假定如此，这个态射的定义并不平凡，因为局部叠 $\XX_{G,F_v}$ 是用 $(\varphi,\Gamma)$-模而非直接用局部伽罗瓦表示定义的。不过，存在一个自然的限制态射，到 Wang-Erickson 的 $p$-进局部叠 [WE18] 的导出版本，而 [BIP23a] 的结果表明这些叠实际上是经典的（比较注 4.1.7）；由注 4.1.2，这些经典叠是我们的局部叠 $\XX_{G,F_v}$ 的子叠。

原则上，我们还应当考虑 $F$ 的实阿基米德素位 $v$ 处的局部叠 $\XX_{G,F_v}$。然而，由于我们假定 $p>2$，这些只是经典叠 [Zhu25a, 命题 2.3.2]，它们分类 ${}^cG$ 中对合的共轭类，这些对合在 ${}^cG$ 到其伽罗瓦因子的投影下位于复共轭 $c_v$ 之上。事实上，我们只想考虑一个特定的共轭类，即 ${}^cG$ 中的「奇」类——见 [BV13, 命题 6.1]。以 $\XX^{\mathrm{odd}}_{G,F,S}$ 记 $\XX_{G,F,S}$ 中分类「在 $F$ 的每个实素位处都为奇」的伽罗瓦表示的那些连通分支的并。

\begin{rem}{9.2.2}\label{rem:a4-9-2-2}
正如 [EZ] 中将会解释的，标准的 Fontaine--Mazur 型猜想蕴含态射 $f$ 可由代数（导出）叠表出。
\end{rem}

\begin{rem}{9.2.3}\label{rem:a4-9-2-3}
$f$ 的切复形（即 $\XX^{\mathrm{odd}}_{G,F,S}$ 在 $\prod_{v\in S}\XX_{G,F_v}$ 上的相对切复形）在某点处（至多差一个平移）由整体到局部伽罗瓦上同调的限制映射的锥来计算；确切地说，我们有一个特异三角（distinguished triangle）
\[
\begin{aligned}
\cdots\to\mathbf t_f&\to C^\bullet(\Gal_{F,S},\mathfrak g)[1]\\
&\to\prod_{v\in S}C^\bullet(\Gamma_{F_v},\mathfrak g)[1]\to\cdots,
\end{aligned}
\]
\end{rem}

\ysrc{344}
于是相对切复形 $\mathbf t_f$ 是 [GV18, 附录 B] 所考虑的那种 Selmer 型复形，它可以有非零的 $H^3$ 项。这个 $H^3$ 项（由整体对偶性）又对偶于另一个伽罗瓦 $H^0$，即 $H^0(\Gal_{F,S},\hat{\mathfrak g}(1))$。（对 $f$ 的这一局部分析，是把整体形变环或提升环表成局部提升环张量积之上的传统方法的一个更几何化的版本。用更传统的提升环语言所作的类似分析见 [Kis07b, §4]。这样的表示在模性提升定理（及其对潜在模性的应用）中很重要；在我们先前的表述中，一个例子是态射 (3.1.3)。）

$f$ 在 $\overline{\rho}$ 处相对拟光滑这一条件很有意思。例如，若 $G=\mathrm{GL}_d$，从而 $\hat{\mathfrak g}=\End(\overline{\rho})$，则 $H^0(\Gal_{F,S},\hat{\mathfrak g}(1))=\End_{\Gal_{F,S}}(\overline{\rho},\overline{\rho}(1))$，于是我们在考虑 $\overline{\rho}$ 是否容许到其 Tate 扭 $\overline{\rho}(1)$ 的非平凡 $\Gal_{F,S}$-等变映射。若 $F=\Q$、$G=\mathrm{GL}_2$，则（回忆 $p>2$）除非 $\overline{\rho}^{\mathrm{ss}}$ 是直和 $1\oplus\overline{\chi}$ 的扭（其中 $\overline{\chi}$ 是 mod $p$ 分圆特征），或者 $p=3$ 且 $\overline{\rho}$ 由 $\Gal_{\Q(\sqrt{-3})}$ 诱导，否则这不可能。这些恰是 Diamond--Taylor [DT94, 定理 1] 关于 $\overline{\rho}$ 提升的局部--整体原理所排除的情形，并且已知此时这种局部--整体原理失效。当 $\overline{\rho}$ 由 $\Gal_{\Q(\sqrt{-3})}$ 诱导时，这要追溯到 Carayol 与 Serre 的工作，[All19, §6.4] 中给出了一个显例。在可约情形，可能层次的问题已由 Ribet 与 Yoo 研究，特别见 [Yoo19]。（$p=2$ 情形的类似结果见 [LHL16]。）

更一般地，$H^0(\Gal_{F,S},\hat{\mathfrak g}(1))=0$ 这一条件与 Taylor--Wiles 方法中的「大像」条件密切相关，也与 [CS17] 的一般性条件密切相关。

\subsection{经由范畴 Langlands 计算志村簇的上同调}
我们陈述一个关于志村簇上同调的猜想性公式，它与 [Zhu25a] §4.7 的诸猜想有关。（该猜想将在 [EZ] 中更详细地讨论。）一般猜想看来远不可及（因为例如它蕴含各种自守性定理！），但我们会在 §9.4 中讨论模曲线的情形。

\subsubsection*{9.3.1.\quad 志村簇的上同调}
现在设 $F=\Q$，并设约化群 $G$ 容许一个 Shimura 数据。我们以 $F$（这个记号现在腾出来了！）记该 Shimura 数据的反射域，并以 $\mu:(\GG_m)_{/F}\to G_{/F}$ 记相伴的 Hodge 余特征标（在共轭意义下良定义）。在 ${}^cG$ 的定义中，我们把 $\Q$ 的扩张 $\tilde F$ 取得足够大，使其包含 $F$，然后写
\[
{}^cG_F:=\hat G\rtimes\bigl(\GG_m\times\Gal(\tilde F/F)\bigr)\subseteq\hat G\rtimes\bigl(\GG_m\times\Gal(\tilde F/\Q)\bigr)=:{}^cG.
\]
把 Hodge 余特征标 $\mu$ 看作 $\hat G$ 的一个权，它给出 $\hat G$ 的一个不可约（最高权）表示 $V_\mu$，后者自然地延拓为 ${}^cG_F$ 的表示。对任何素位集 $S$（照例包含 $p$ 与 $\infty$），模叠 $\XX_{G,\Q,S}$ 参数化（导出的）C-同态 $\Gal_\Q\to{}^cG$ 的一个万有族，我们可以把它限制到 $\Gal_F$，从而得到一族表示 $\Gal_F\to{}^cG_F$。把这族表示与 ${}^cG_F$ 的表示 $V_\mu$ 复合，便得到 $\XX_{G,\Q,S}$ 上的向量丛 $\mathcal V_\mu$，它带有一个 $\Gal_F$-作用。

\ysrc{345}
选取一个级子群（level subgroup）$K_f$，即 $G(\mathbf A)$ 的紧开子群\tn{底本此处作 $G(\mathbf A)$，但下文 $K_f$ 只对有限素位作分解，且记号本身表明它应在 $G(\mathbf A_f)$ 中；疑为 $G(\mathbf A_f)$ 之误，译文照原（上下文推断，无外部依据）。}，我们假定它可以分解为 $K_f=\prod_vK_v$，其中对每个有限素位 $v$，$K_v$ 是 $G(\Q_v)$ 的紧开子群。我们（尽其所可）把 $S$ 取得足够大，使 $G$ 在 $S$ 外非分歧，并且存在 $G$ 到 $\Z[1/S]$ 上的约化群概形的扩张 $\mathcal G$，使得对 $v\notin S$ 有 $K_v=\mathcal G(\Z_v)$。局部对称商 $Y(K_f):=G(\Q)\backslash\bigl(X\times G(\mathbf A_f)\bigr)/K_f$ 是一个志村簇。

对每个有限素位 $v\in S$，设 $W_v$ 是 $K_v$ 在有限挠 $\OO$-模上的光滑表示。以 $W:=\bigotimes_{v\in S}W_v$ 记 $K_f$ 的相应表示（它经由到 $\prod_{v\in S}K_v$ 的投影作用；当然，张量积在 $\OO$ 上取），并以 $\mathcal W$ 记 $Y(K_f)$ 上相伴的局部系统 $\mathcal W:=G(\Q)\backslash\bigl(X\times G(\mathbf A_f)\times W\bigr)/K_f$。对每个 $v$，紧诱导 $c\text{-}\Ind_{K_v}^{G(\Q_v)}W_v$ 是光滑的 $G(\Q_v)$-表示。施加猜想 6.1.15 与 8.3.1 中的函子，我们（猜想性地！）得到 $\XX_{G,\Q_v}$ 上的一个凝聚层，就简单地记作 $\mathfrak A_v$。我们可以构造外张量积 $\mathfrak A:=\boxtimes_{v\in S_f}\mathfrak A_v$，它是 $\prod_{v\in S_f}\XX_{G,\Q_v}$ 上的对象。

如上所述，以 $f:\XX^{\mathrm{odd}}_{G,\Q,S}\to\prod_{v\in S}\XX_{G,\Q_v}$ 记整体到局部的限制态射。于是我们对志村簇的上同调有如下猜想性公式。

\begin{conj}{9.3.2}\label{conj:a4-9-3-2}
设 $d$ 是 $Y(K_f)$ 的复维数。则存在同构
\begin{equation}\tag{9.3.3}\label{eq:a4-9-3-3}
R\Gamma_c\bigl(Y(K_f),\mathcal W\bigr)[d]\isom R\Gamma\bigl(\XX^{\mathrm{odd}}_{G,\Q,S},f^!\mathfrak A\otimes\mathcal V_\mu\bigr),
\end{equation}
它与源与靶上的 $\Gal_F$-作用及（导出）Hecke 作用相容。
\end{conj}

通过取 mod $p^n$ 情形的逆向极限，该猜想自然地延拓到 $p$-进系数的上同调情形。我们还可以对 $S$ 中素位处的驯顺层次取极限，以研究完备上同调，如下文的预期定理 9.4.2 那样。

\begin{rem}{9.3.4}\label{rem:a4-9-3-4}
正如我们在 §8.3 讨论猜想 8.3.1 时简要提到的，猜想 9.3.2 陈述中出现的局部函子依赖于 Whittaker 型辅助数据的选择。在猜想的陈述中，我们应当对这些数据作一个整体选择，然后用所诱导的局部数据计算局部函子 $\mathfrak A_v$。所得的层 $f^!\mathfrak A$ 于是应当良定义，不依赖于整体的选择。（整体数据的不同选择应当使每个 $\mathfrak A_v$ 的取值被某个线丛作扭，但这些不同线丛拉到 $\XX^{\mathrm{odd}}_{G,\Q,S}$ 上的张量积应当是平凡的；函数域语境中的类似讨论见 [Zhu25a]，猜想 4.7.5 之前的讨论。）
\end{rem}

\begin{rem}{9.3.5}\label{rem:a4-9-3-5}
猜想 9.3.2 特别蕴含 $\Gal_F$ 的种种表示出现在志村簇的上同调中，从而特别蕴含 Fontaine--Mazur 猜想的某些情形。要使这一蕴含精确化，需要知道关于层 $\mathfrak A_v$ 的支撑的某些信息。对 $v\mid p$，猜想 6.1.15 (4) 是相关的陈述：它表明 $\mathfrak A_v$ 会探测 $\XX^{\mathrm{odd}}_{G,\Q,S}$ 中那些到 $\XX_v$ 的限制满足适当 $p$-进 Hodge 理论条件的点。$v\nmid p$ 情形的类似陈述（用分歧条件描述 $\mathfrak A_v$ 的支撑）也应当成立。我们不在此写出这些陈述，而请读者参看 [BZCHN24]、[FS24]、\ysrc{346}[Hel23] 与 [Zhu25a] 中关于 $v\nmid p$ 情形的讨论；此时的关键点在于，范畴理论应当在适当的意义下与经典局部 Langlands 对应相容。
\end{rem}

\begin{rem}{9.3.6}\label{rem:a4-9-3-6}
自然地作用在 (9.3.3) 的源与靶上的 Hecke 代数，是对有限素位 $v$ 的导出 Hecke 代数 $\mathrm{RHom}_G\bigl(c\text{-}\Ind_{K_v}^{G(\Q_v)}W_v,c\text{-}\Ind_{K_v}^{G(\Q_v)}W_v\bigr)$（其中当 $v\notin S$ 时，我们简单地把 $W_v$ 解释为 $K_v$ 的平凡表示；要看出导出 Hecke 代数确实在 $v\notin S$ 处作用，可以比较 [Zhu25a]，猜想 4.7.5 之后紧接的 $S$ 与 $S\cup\{v\}$ 两种情形）。关于这些导出 Hecke 代数以及与之相关的若干猜想（在 $v\neq p$ 的情形），请读者参看 [Zhu25a]。
\end{rem}

\begin{rem}{9.3.7}\label{rem:a4-9-3-7}
因子 $\mathcal V_\mu$ 带有整体的一面，因为它带有整体绝对伽罗瓦群 $\Gal_F$ 的作用；但若忽略这一点，它就可以看作 $\infty$ 处对猜想 9.3.2 公式的贡献。我们猜测，对任何连通约化群 $G$（即对一切局部对称空间 $Y(K_f)$ 的上同调，无论它们是否具有志村簇的结构），只要适当选取来自 $\infty$ 的因子（该因子此时大概只是 $\XX^{\mathrm{odd}}_{G,\Q,S}$ 上的一个向量丛，由 $\infty$ 处奇模叠上适当选取的向量丛拉回而得，并且不再带有伽罗瓦表示的额外结构），就应当有类似的同构。特别地，我们预期可以重新得到 Galatius 与 Venkatesh 在 [GV18, Ven19] 中观察并猜想的那些导出现象。
\end{rem}

\subsection{整体情形的一个例子：$(\mathrm{PGL}_2)_{/\Q}$}
在 [EGZ] 中，M.E.、T.G. 与 Xinwen Zhu 用 Taylor--Wiles--Kisin 拼接方法研究模曲线情形的猜想 9.3.2。我们在一个特殊情形下简要解释这项工作的一部分（它仍在进行中），它与我们在 §3 中对拼接方法的讨论相适配；另见 [JNWE24, §6]。

\subsubsection*{9.4.1.\quad 与拼接的关系}
为明确起见，置 $G=\mathrm{PGL}_2$，于是我们处于模曲线的情形。按 §3.1 选取 $\overline{r}:\Gal_{\Q,\{p,\infty\}}\to\mathrm{GL}_2(k)$，特别地 $\overline{r}|_{\Gal_{\Q(\zeta_p)}}$ 绝对不可约，且 $\overline{r}$ 的行列式为 $\overline{\varepsilon}^{-1}$。记 $\overline{\rho}:=\overline{r}|_{\Gal_{\Q_p}}$。我们考虑 (9.2.1) 在 $S=\{p,\infty\}$ 时的特殊情形，于是有态射
\[
f:\XX^{\varepsilon^{-1},\mathrm{odd}}_{\mathrm{GL}_2,\Q,\{p,\infty\}}\to\XX^{\varepsilon^{-1}}_{\mathrm{GL}_2,\Q_p},
\]
其中我们只考虑行列式为 $\varepsilon^{-1}$ 的表示，因而目标是 §7.2 中记作 $\XX$ 的那个叠。以 $\mathcal V$ 记 $\XX^{\varepsilon^{-1},\mathrm{odd}}_{\mathrm{GL}_2,\Q,\{p,\infty\}}$ 上由万有 $\Gal_\Q$-表示给出的秩 $2$ 向量丛。

如同 §3.1，以 $R(\overline{\rho})$ 记 $\overline{\rho}$ 的固定行列式 $\varepsilon^{-1}$ 的万有形变环。按定义存在 versal 态射 $g:\Spf R(\overline{r}_p)\to\XX_p$\tn{底本此处先作 $R(\overline{r}_p)$、后作 $R(\overline{r}|_{\Gal_{\Q_p}})$，而据上文 $\overline{\rho}=\overline{r}|_{\Gal_{\Q_p}}$，两处当指同一个环；译文照录原文（上下文推断，无外部依据）。}，它把其定义域中的闭点映到局部限制 $f(\overline{r})$。拼接构造给出一个 $R(\overline{r}|_{\Gal_{\Q_p}})$-模 $M_\infty$，带有 $\mathrm{PGL}_2(\Q_p)$ 的作用。

我们预期在 [EGZ] 中证明下述结果，它描述完备同调群 $\widehat H_1$ 在与 $\overline{r}$ 对应的、与 $p$ 互素的 Hecke 代数的极大理想 $\mathfrak m$ 处的局部化。如同 (9.1.1)，我们以 $\XX_{\overline{r}}$ 记 $\XX^{\varepsilon^{-1},\mathrm{odd}}_{\mathrm{GL}_2,\Q,\{p,\infty\}}$ 中与 $\overline{r}$ 对应的不可约分支。

% 本片止于印刷页 346 页末（§9.4.1 中段），续片自 PDF 360 页起
